\scanpage{16}
ne peut suffire à reconstruire $Q$ à isom. unique près -- mais seulement $Q$ ``modulo antipodisme''. Considérons le carré épinglé $Q_0$, et le groupe de ses automorphismes
\[
\deleted{\mathfrak{D}_4}\quad D_4=\operatorname{Aut}(Q_0)
\]
avec la suite exacte canonique
\diagram{\begin{tikzcd}
1 \arrow[r] & \mathbb{F}_2 \arrow[r] & D_4 \arrow[r] & \mathbb{F}_2^2 \arrow[r,dash] & 1
\end{tikzcd}}
et la donnée de notre $\mathfrak{D}$ revient à celle d'un torseur sous $\mathbb{F}_2^2$ \inserted{le quotient de $D_4$} (i.e. de deux ens. à deux éléments (ce sont en fait $D$ et $C$, dans cet ordre\ldots)), tandis que celle d'un carré revient à la donnée d'un torseur sous $D_4$ lui-$=$.

On a une opération \inserted{(par autom. int.)} de $\mathfrak{D}^+$ sur $D$\inserted{$\hookrightarrow\mathfrak{D}^+$}, la donnée de $Q$ (correspondant au groupe d'autom $\mathfrak{D}$) \deleted{\unsure{nous}} revient à celle d'un relèvement de \deleted{\unsure{$\mathfrak{D}$}} l'opération \deleted{\unsure{induite de}} en \deleted{\unsure{les}} une opération \deleted{\unsure{sur}} transitive sur un ens. $S$, \inserted{(du \unsure{m.|même} gr. bien sûr $\mathfrak{D}$)} avec $S/\underline{a}\simeq D$.
Description duale via $C\hookrightarrow\mathfrak{D}^+$\ldots

\uline{Remarque.} Le groupe $\operatorname{Aut}(D_4)$ est isom. à $D_4$ (mais \uline{pas} via autom. intérieurs, bien sûr), ce qui implique qu'il y a \inserted{quand même} une équivalence de groupoïdes
\[
(\text{Carrés})\xrightarrow{\approx}(\text{groupes isom. à }D_4)
\]
(mais qui n'est \unsure{$\not=$} donnée par \deleted{\unsure{$\operatorname{Aut}_Q$}} $Q\mapsto\operatorname{Aut}(Q)$ !), qui mériterait d'être comprise\ldots
% archive foliation: 15
